% ============================================================
%  Kapitel 6: Zusammenfassung
% ============================================================
\chapter{Zusammenfassung}

% ============================================================
\section{Architektur -- Rueckblick}
% ============================================================

TilVista v0.5.42 demonstriert eine durchgaengige,
signalgesteuerte Qt6-Architektur. Die wichtigsten
Entwurfsentscheidungen im Rueckblick:

\subsection{Entkopplung durch std::function-Callbacks}

Kein Tab importiert den Header eines anderen Tabs.
Statt \verb|MainWindow*| als Parameter erhalten alle
Tabs ein Lambda:

\begin{lstlisting}
// In MainWindow::MainWindow():
m_aleaVueTab = new AleaVueTab(
    [this]{ return m_dirBar->directory(); },
    m_shujuko);
\end{lstlisting}

Das entspricht dem \emph{Dependency Inversion Principle}
(DIP) aus den SOLID-Prinzipien: Abhaengigkeit von einer
Abstraktion (Funktion), nicht von einer konkreten Klasse.

\subsection{Unidirektionale Signalkette}

Zustandsaenderungen fliessen immer von einem Sender
durch \klasse{MainWindow} zu den betroffenen Empfaengern
-- niemals direkt zwischen Tabs:

\begin{center}
\small
\begin{tabular}{ccc}
\texttt{DirDatabasePanel} &
\(\longrightarrow\) \textit{Signal} \(\longrightarrow\) &
\texttt{MainWindow} \\
& & \(\downarrow\) \\
& & \texttt{SattumaPicTab} \\
& & \texttt{MadoludusTab} \\
\end{tabular}
\end{center}

Kein Tab reagiert direkt auf Signale eines anderen Tabs.
Das macht den Datenfluss nachvollziehbar und
Fehlerquellen leicht lokalisierbar.

\subsection{Worker-Thread-Muster durchgaengig}

Jede datei-IO-Operation laeuft in einem eigenen
\klasse{QThread} mit dem Worker-Object-Muster:

\begin{lstlisting}
auto* worker = new MeinWorker(parameter);
auto* thread = new QThread;
worker->moveToThread(thread);
connect(thread, &QThread::started,
        worker, &MeinWorker::run);
connect(worker, &MeinWorker::resultReady,
        this,   &MeineKlasse::onErgebnis);
connect(worker, &MeinWorker::resultReady,
        thread, &QThread::quit);
connect(worker, &MeinWorker::resultReady,
        worker, &QObject::deleteLater);
connect(thread, &QThread::finished,
        thread, &QObject::deleteLater);
thread->start();
\end{lstlisting}

Der GUI-Thread blockiert nie. Fortschritt kommt per
\signal{progressChanged(int)}, Ergebnis per
\signal{resultReady}.

\subsection{safeStop-Idiom}

Vor jedem Thread-Start werden alle Signal-Verbindungen
des alten Threads zum Empfaenger getrennt:

\begin{lstlisting}
static void safeStop(QThread*& t, QObject* receiver) {
    if (!t) return;
    t->disconnect(receiver);
    t = nullptr;
}
\end{lstlisting}

Das verhindert Stale-Signal-Bugs ohne Mutex oder
\verb|QThread::wait()| (das wegen Queued-Connections
deadlocken wuerde).

\subsection{Optionale Abhaengigkeit: TV\_NO\_MULTIMEDIA}

Qt Multimedia wird per \cmake{find\_package QUIET}
gesucht. Nicht gefunden: \makro{TV\_NO\_MULTIMEDIA}
wird als Praeprozessor-Makro definiert. Im Code
steuern \verb|#ifndef|/\verb|#ifdef| die Codepfade.
Das Programm kompiliert und laeuft in beiden Faellen.

\subsection{v0.5.42: MadoLudus-Refactor}

Die Trennung von \klasse{MadoludusTab}
(Konfiguration) und \klasse{MadoludusWindow}
(Wiedergabe) reduzierte eine 800-Zeilen-Monolith-Klasse
auf zwei klar abgegrenzte Verantwortlichkeiten.
Das Fenster kommuniziert per zwei Signalen zurueck:
\signal{windowClosed()} und
\signal{currentFileChanged(QString)}.

% ============================================================
\section{C++-Konzepte im Ueberblick}
% ============================================================

\begin{longtable}{p{4.2cm}p{4.2cm}p{5cm}}
\toprule
\textbf{Konzept} & \textbf{Datei(en)} & \textbf{Zweck} \\
\midrule
\endfirsthead
\midrule
\endhead
\verb|std::function<T>| &
  aleavuetab, sattumapictab, madoludustab &
  Callback ohne harte Abhaengigkeit \\[4pt]

\verb|std::mt19937| + \verb|std::random_device| &
  sattumapictab, slideshowwindow, madoluduswindow &
  Qualitaetszufall, OS-Entropie als Seed \\[4pt]

\verb|static| (lokal) &
  pathutils, DirDatabasePanel &
  Einmalige Initialisierung, Singleton-Muster \\[4pt]

\verb|static constexpr| &
  slideshowwindow, videopreviewwidget, madoluduswindow &
  Compile-Zeit-Konstanten, typsicher \\[4pt]

\verb|= nullptr| / \verb|= false| (In-class-Init) &
  alle Klassen &
  Sichere Member-Initialisierung ohne Konstruktor \\[4pt]

\verb|override| &
  alle Qt-Subklassen &
  Compiler prueft korrekte Methodensignatur \\[4pt]

Raw-String \verb|R"(...)"| &
  pathutils, videoframesworker &
  Sonderzeichen ohne Escape-Sequenzen \\[4pt]

\verb|enum class| &
  videoframesworker &
  Stark typisiertes Enum, kein implizites \verb|int| \\[4pt]

\verb|try/catch(...)| &
  scanworker &
  Alle Exceptions im Worker-Thread abfangen \\[4pt]

Template-Funktion &
  dirdatabasepanel &
  DRY: generischer Thread-Starter \verb|launchWorker<>| \\[4pt]

\verb|sizeof(arr)/sizeof(arr[0])| &
  madoludustab &
  Compile-Zeit-Arraylaenge \\[4pt]

\verb|explicit| Konstruktor &
  alle Klassen &
  Verhindert ungewollte implizite Konvertierungen \\[4pt]

\verb|QStringLiteral| &
  pathutils &
  Compile-Zeit-String, schneller als Laufzeit-Konstruktion \\[4pt]

\verb|struct| mit Defaults &
  madoluduswindow &
  Konfigurationsdaten ohne Konstruktor-Boilerplate \\[4pt]

\bottomrule
\end{longtable}

% ============================================================
\section{Qt-Konzepte im Ueberblick}
% ============================================================

\begin{longtable}{p{4.0cm}p{4.4cm}p{5cm}}
\toprule
\textbf{Klasse / Konzept} & \textbf{Datei(en)} & \textbf{Zweck} \\
\midrule
\endhead

\klasse{QMainWindow} &
  mainwindow, slideshowwindow, madoluduswindow &
  Vollstaendiges Fenster mit Status\-leiste \\[4pt]

\klasse{QWidget} &
  alle Tab-Klassen, Panels &
  Basis aller GUI-Elemente \\[4pt]

\klasse{QTabWidget} &
  mainwindow &
  Tab-Navigation zwischen den Werkzeugen \\[4pt]

\klasse{QStackedWidget} &
  videopreviewwidget &
  Zwischen Seiten umschalten (Video / Bild) \\[4pt]

\klasse{QTimer} &
  slideshowwindow, madoluduswindow, videopreview &
  Zeitgesteuerte Ereignisse ohne Blocker \\[4pt]

\klasse{QThread} + moveToThread &
  alle Worker &
  Hintergrundarbeit ohne GUI-Blockade \\[4pt]

Queued Connection &
  alle Worker-Signale &
  Automatisch thread-sicherer Signal-Transport \\[4pt]

\verb|deleteLater()| &
  alle Worker &
  Thread-sicheres Aufraeuimen von QObjects \\[4pt]

\klasse{QDirIterator} &
  scanworker &
  Rekursive Verzeichnis-Traversal \\[4pt]

\klasse{QImageReader} &
  slideshowwindow, sattumapictab, madoluduswindow &
  Effizientes Laden mit Vorab-Skalierung \\[4pt]

\klasse{QPainter} &
  madooverlay, previewlabel &
  2D-Zeichenoperationen, Custom-Painting \\[4pt]

\klasse{QJsonObject} / \klasse{QJsonDocument} &
  config, dbworkers, shujukopanel &
  JSON-Persistenz, Copy-on-Write \\[4pt]

\klasse{QProcess} &
  videoframesworker &
  Externe Programme (ffmpeg) starten \\[4pt]

\klasse{QRegularExpression} &
  videoframesworker, dirdatabasepanel &
  PCRE2-Regex, statisch kompiliert \\[4pt]

\klasse{QShortcut} &
  mainwindow &
  App-weite Tastenkuerzel \\[4pt]

\klasse{QFileIconProvider} &
  sattumapictab &
  Native OS-Datei-Icons \\[4pt]

\klasse{QMediaPlayer} + \klasse{QVideoWidget} &
  videopreviewwidget, madoluduswindow &
  Native Videowiedergabe \\[4pt]

\makro{Q\_OS\_WIN} / \makro{TV\_NO\_MULTIMEDIA} &
  pathutils, alle Multimedia-Klassen &
  Plattform- und Feature-bedingte Compilation \\[4pt]

\verb|Qt::ApplicationShortcut| &
  mainwindow &
  Kuerzel auch bei Fokus in Eingabefeld aktiv \\[4pt]

\verb|Qt::FramelessWindowHint| &
  madoluduswindow &
  Rahmenlos + manuell verschiebbar \\[4pt]

\verb|WA_DeleteOnClose| &
  madoluduswindow, slideshowwindow &
  Automatisches Loeschen beim Schliessen \\[4pt]

\bottomrule
\end{longtable}

% ============================================================
\section{Empfohlener Lernpfad}
% ============================================================

Fuer den Einstieg in TilVista als Lernprojekt empfiehlt
sich folgende Lesereihenfolge -- von einfach nach komplex:

\begin{enumerate}
  \item \textbf{main.cpp} --
    Einstiegspunkt, \klasse{QApplication}, \klasse{Config},
    erstes Fenster; zwei Dutzend Zeilen die das grosse Bild
    zeigen

  \item \textbf{pathutils.h + pathutils.cpp} --
    Namespace, statische lokale Variablen,
    \verb|#ifdef Q_OS_WIN|, Raw-Strings; keine Qt-Abhaengigkeit
    ausser \klasse{QString}

  \item \textbf{config.h + config.cpp} --
    Statische Membervariablen, \klasse{QJsonDocument},
    einfachstes moegliches Singleton

  \item \textbf{scanworker.h + scanworker.cpp} --
    Erstes vollstaendiges Worker-Thread-Beispiel,
    \klasse{QDirIterator}, Signal-Fortschritt,
    \verb|try/catch|

  \item \textbf{mainwindow.h + mainwindow.cpp} --
    \klasse{QMainWindow}, Layout, Lambda-Callbacks,
    Signalkette, \verb|closeEvent|

  \item \textbf{slideshowwindow.h + slideshowwindow.cpp} --
    \klasse{QTimer}, \verb|keyPressEvent|,
    \klasse{QImageReader} mit Vorab-Skalierung,
    Verlaufslogik, Vollbild-Cursor

  \item \textbf{dirdatabasepanel.h + dirdatabasepanel.cpp} --
    JSON-Persistenz, drei parallele Threads,
    \verb|safeStop|, \verb|setBusy|,
    \verb|flushPendingWrites|, Template-Helfer

  \item \textbf{shujukopanel.h + shujukopanel.cpp} --
    Zweites JSON-Schema, Bool-Guard,
    \verb|validateFiles|, \verb|flushPendingWrites|

  \item \textbf{videoframesworker.h + videoframesworker.cpp} --
    \klasse{QProcess}, \verb|enum class|,
    dauer-bewusste Timestamps, Regex auf stderr

  \item \textbf{madoluduswindow.h + madoluduswindow.cpp} --
    Rahmenlos, Drag-to-Move, History-Navigation,
    optionale Multimedia, \verb|WA\_DeleteOnClose|

  \item \textbf{madoludustab.h + madoludustab.cpp} --
    Konfiguration trennen von Ausfuehrung,
    \verb|buildConfig()|, \verb|sizeof|-Trick

  \item \textbf{madooverlay.h + madooverlay.cpp} --
    Custom \verb|paintEvent|, Unicode-Emojis,
    Stylesheet-Zustandsaenderung per Methode
\end{enumerate}

% ============================================================
\section{Weiterfuehrende Ressourcen}
% ============================================================

\begin{description}[leftmargin=3cm, style=nextline]
  \item[Qt-Dokumentation]
    \url{https://doc.qt.io/qt-6/}\\
    Vollstaendige API-Referenz, Beispiele und Guides

  \item[Signals \& Slots]
    \url{https://doc.qt.io/qt-6/signalsandslots.html}\\
    Ausfuehrliche Erklaerung beider Syntaxen

  \item[Threading-Guide]
    \url{https://doc.qt.io/qt-6/thread-basics.html}\\
    Worker-Object-Muster, QThread-Lebenszyklus

  \item[Model/View]
    \url{https://doc.qt.io/qt-6/model-view-programming.html}\\
    Naechster Schritt nach dem Widget-basierten Ansatz

  \item[cppreference C++17]
    \url{https://en.cppreference.com/}\\
    Lambda, std::function, constexpr, enum class

  \item[C++ Core Guidelines]
    \url{https://isocpp.github.io/CppCoreGuidelines/}\\
    Empfohlene Idiome (RAII, Rule of Five, etc.)

  \item[CMake-Dokumentation]
    \url{https://cmake.org/cmake/help/latest/}\\
    find\_package, Generator-Expressions, Funktionen

  \item[CTAN: babel-german]
    \url{https://ctan.org/pkg/babel-german}\\
    Installation des fehlenden deutschen Babel-Pakets
    (dann kann in \datei{packages.tex} \verb|\usepackage[ngerman]{babel}|
    aktiviert werden)
\end{description}
